\begin{abstract}
Vision-Language-Action (VLA) models report high Collision Avoidance Rates (CAR) on safety benchmarks, yet we show these metrics are contaminated by action collapse, where near-zero motor output produces zero collisions through inaction rather than avoidance.
On SafeLIBERO, an SDF-augmented Pi0 achieves 91.5\% CAR while 99.6\% of its safe episodes exhibit end-effector (EE) displacement below 0.1\,mm. The robot is action-collapsed, not safe.
We introduce displacement-filtered CAR (dfCAR), which excludes sub-threshold episodes. The bimodal displacement distribution makes dfCAR robust across a $[2, 200]$\,mm threshold range.
We further show that end-effector-only Signed Distance Field (SDF) supervision fails to deliver collision avoidance. At inference time, ground-truth oracle refinement reduces CAR below unrefined baselines because the SDF gradient drives unmonitored arm links into obstacles. During training, collision-free demonstrations alone account for nearly all measured avoidance gains (${\sim}$26 percentage points). SDF supervision shows no statistically significant benefit at the tested scale $\lambda \in \{0.005, 0.01, 0.1\}$ ($p > 0.05$ for $\lambda \in \{0.01, 0.1\}$; $\lambda = 0.1$: 57.00\%$\,\pm\,$2.18, 3 seeds).
Displacement filtering is necessary for meaningful VLA safety evaluation. EE-only geometric supervision does not significantly improve collision avoidance across the tested loss-weight range, suggesting that collision avoidance primarily arises from demonstration data rather than EE-only geometric supervision.
\end{abstract}
